\subsection{模型检验与尾部风险}

短期检验方面，模型同时通过基准对比、滞后平移、机制消融和参数扰动复核。局部参数扰动表明，最优参数附近预测路径仍围绕真实价格平台展开，说明结论不依赖某一个极窄参数点。R 语言独立计量检验进一步复核了主模型相对朴素基准的误差优势和残差结构。

长期检验方面，蒙特卡洛情景树给出第 180 天价格 5\%--95\% 区间 87.77--116.59 美元/桶，外推期最高价突破 120 美元/桶的条件概率约为 20.0\%。状态转移情景树进一步引入缓和、维持、升级三类状态，第 180 天价格 5\%--95\% 区间为 90.98--116.21 美元/桶，外推期最高价 P95 约为 137.62 美元/桶，突破 120 美元/桶的条件概率约为 85.4\%。

\begin{table}[!htbp]
\centering
\caption{模型检验与可靠性摘要}
\begin{tabularx}{\linewidth}{lX}
\toprule
检验维度 & 主要结论 \\
\midrule
基准对比 & 短期冲击模型 RMSE 低于朴素上一日基准和滚动 ARIMA 基准。 \\
滞后检验 & 原始位置误差最低，预测曲线不存在明显滞后复制。 \\
机制消融 & 关闭风险、缓冲确认或预期修复机制会显著改变误差结构。 \\
参数扰动 & 局部扰动下价格平台仍稳定，说明结论不依赖单点参数。 \\
尾部风险 & 长期结果需结合三情景、蒙特卡洛分位区间和状态转移概率理解。 \\
外生变量 & GPR 与 OVX 用于风险环境和长期状态约束，不替代附件真实价格。 \\
\bottomrule
\end{tabularx}
\end{table}

\keypoint{短期模型给出可复核的事件窗口拟合，长期模型给出条件情景和尾部风险边界。模型结论应读作霍尔木兹封锁条件下的机制解释和风险评估，而非跨年份通用油价交易预测。}
